% This is text associated with the open textbook "Open Electromagnetics"
% See immediately below for author, date
% This work is licensed under the Creative Commons Attribution-ShareAlike 4.0 International License. (CC BY SA 4.0) 
% To view a copy of the license, visit http://creativecommons.org/licenses/by-sa/4.0/

%ORB TITLE Permittivity of Some Common Materials
%ORB AUTHOR 0        # S.W. Ellingson, Virginia Tech (ellingson@vt.edu)
%ORB DATE 20191128   # yyyymmdd
%ORB ABSTRACT Permittivity of Some Commonly-Encountered Materials Materials

\label{m0135_Table_of_Permittivities} % <-- Must match name of this file and module directory name.  Don't mess with this!
\index{permittivity!of common materials|(}

The values below are relative permittivity $\epsilon_r \triangleq \epsilon/\epsilon_0$ for a few materials that are commonly encountered in electrical engineering applications, and for which permittivity emerges as a consideration.
Note that ``relative permittivity'' 
\index{permittivity!relative}
is sometimes referred to as \emph{dielectric constant}.
\index{dielectric constant}

Here we consider only the physical (real-valued) permittivity, which is the real part of the complex permittivity (typically indicated as $\epsilon'$ or $\epsilon_r'$) for materials exhibiting significant loss.

Permittivity varies significantly as a function of frequency.
The values below are representative of frequencies from a few kHz to about 1~GHz.
The values given are also representative of optical frequencies for materials such as silica that are used in optical applications.
Permittivity also varies as a function of temperature.
In applications where precision better than about 10\% is required, primary references accounting for frequency and temperature should be consulted.
The values presented here are gathered from a variety of references, including those indicated in ``Additional References.''
  
{\bf Free Space} (vacuum): $\epsilon_r \triangleq 1$

\vfill\break
{\bf Solid Dielectrics}:
\index{dielectric (material)!examples}
%
\begin{tabular}{lll}
Material                   & $\epsilon_r$ & Common uses \\
\hline
Styrofoam$^1$              & 1.1 & \\ \index{styrofoam} % https://en.wikipedia.org/wiki/Styrofoam 
Teflon$^2$                 & 2.1 & \\ \index{Teflon} % HPC, https://en.wikipedia.org/wiki/Polytetrafluoroethylene
Polyethylene               & 2.3 & coaxial cable \index{polyethylene} \\ % HPC
Polypropylene              & 2.3 & \\ \index{polypropylene} % HPC
Silica                     & 2.4 & optical fiber$^3$ \index{silica} \\ %https://en.wikipedia.org/wiki/Optical_fiber
Polystyrene                & 2.6 & \\ \index{polystyrene} % HPC, https://en.wikipedia.org/wiki/Polystyrene
Polycarbonate              & 2.8 & \\ \index{polycarbonate} % HPC, https://en.wikipedia.org/wiki/Polycarbonate
Rogers RO3003              & 3.0 & PCB substrate \\ \index{printed circuit board (PCB)!material}
FR4 (glass epoxy laminate) & 4.5 & PCB substrate \index{FR4} \\ % https://en.wikipedia.org/wiki/FR-4
\end{tabular}
%
~\\
$^1$ Properly known as \emph{extruded polystyrene foam} (XPS).\\
$^2$ Properly known as \emph{polytetrafluoroethylene} (PTFE).\\ \index{polytetrafluoroethylene}
$^3$ Typically doped with small amounts of other materials to slightly raise or lower the index of refraction 
\index{index of refraction}
($=\sqrt{\epsilon_r}$).

Non-conducting spacing materials used in discrete capacitors exhibit $\epsilon_r$ ranging from about 5 to 50.

{\bf Semiconductors} 
\index{semiconductor (material)}
commonly appearing in electronics -- including carbon, silicon, 
\index{silicon}
geranium, 
\index{geranium}
indium phosphide, 
\index{indium phosphide}
and so on -- typically exhibit $\epsilon_r$ in the range 5--15.  
% https://www.microwaves101.com/encyclopedias/indium-phosphide

{\bf Glass} 
\index{glass}
exhibits $\epsilon_r$ in the range 4--10, depending on composition.

{\bf Gasses}, including air, typically exhibit $\epsilon_r\cong 1$ to within a tiny fraction of a percent. 

{\bf Liquid water} 
\index{water}
typically exhibits $\epsilon_r$ in the range 72--81. Distilled water exhibits $\epsilon_r \approx 81$ at room temperature, whereas sea water tends to be at the lower end of the range.    

{\bf Other liquids} typically exhibit $\epsilon_r$ in the range 10--90, with considerable variation as a function of temperature and frequency.  Animal flesh and blood consists primarily of liquid matter and so also exhibits permittivity in this range.

{\bf Soil} 
\index{soil}
typically exhibits $\epsilon_r$ in the range 2.5--3.5 when dry and higher when wet.  The permittivity of soil varies considerably depending on composition.

%--------------------------------

{\bf Additional Reading:}
\begin{itemize}
\item \href{http://hbcponline.com}{\emph{CRC Handbook of Chemistry and Physics}}.
\item ``\href{https://en.wikipedia.org/wiki/Relative_permittivity}{Relative permittivity}'' on Wikipedia.
\item ``\href{http://www.microwaves101.com/encyclopedias/miscellaneous-dielectric-constants}{Miscellaneous Dielectric Constants}'' on \href{http://www.microwaves101.com}{microwaves101.com}.
\end{itemize}

\index{permittivity!of common materials|)}


